% ============================================================================
% thesis_template.tex - Minimal Thesis Template (book class)
% ============================================================================
% A minimal but complete thesis template demonstrating:
%   - book document class with frontmatter / mainmatter / backmatter
%   - Chapter-level organisation
%   - Appendix
%   - Bibliography with BibTeX
%
% Compilation:
%   pdflatex thesis_template.tex
%   bibtex   thesis_template
%   pdflatex thesis_template.tex
%   pdflatex thesis_template.tex
%
% (Or: latexmk -pdf thesis_template.tex)
% ============================================================================

\documentclass[
    12pt,
    a4paper,
    twoside,           % mirror margins for binding
    openright          % chapters start on right-hand (odd) pages
]{book}

% ============================================================================
% PACKAGES
% ============================================================================

\usepackage[utf8]{inputenc}
\usepackage[T1]{fontenc}
\usepackage{lmodern}
\usepackage[english]{babel}

% Layout
\usepackage[
    top=2.5cm,
    bottom=2.5cm,
    inner=3.5cm,       % extra space for binding
    outer=2.5cm,
    headheight=14pt
]{geometry}

% Headers and footers
\usepackage{fancyhdr}

% Typography
\usepackage{microtype}
\usepackage{setspace}
\onehalfspacing        % 1.5× line spacing throughout

% Mathematics
\usepackage{amsmath, amssymb, amsthm}

% Figures and graphics
\usepackage{graphicx}
\usepackage{float}
\usepackage{caption}
\captionsetup{font=small, labelfont=bf}

% Tables
\usepackage{booktabs}

% Colors
\usepackage{xcolor}
\definecolor{thesisblue}{RGB}{0, 56, 101}   % Institutional dark blue

% Cross-references
\usepackage[
    hidelinks,
    pdftitle={On the Convergence of Iterative Methods for Sparse Linear Systems},
    pdfauthor={Jane Doe},
    pdfsubject={Mathematics Thesis},
    pdfkeywords={numerical methods, sparse matrices, convergence, LaTeX}
]{hyperref}
\usepackage{cleveref}   % \cref, \Cref for smart cross-references

% ============================================================================
% THEOREM ENVIRONMENTS
% ============================================================================

\newtheorem{theorem}{Theorem}[chapter]
\newtheorem{lemma}[theorem]{Lemma}
\newtheorem{proposition}[theorem]{Proposition}
\newtheorem{corollary}[theorem]{Corollary}
\newtheorem{definition}[theorem]{Definition}
\newtheorem{remark}[theorem]{Remark}
\newtheorem{example}[theorem]{Example}

% ============================================================================
% CUSTOM COMMANDS
% ============================================================================

\newcommand{\R}{\mathbb{R}}           % Real numbers
\newcommand{\N}{\mathbb{N}}           % Natural numbers
\newcommand{\norm}[1]{\lVert #1 \rVert}  % Vector norm
\newcommand{\abs}[1]{\lvert #1 \rvert}   % Absolute value

% ============================================================================
% PAGE STYLES
% ============================================================================

\pagestyle{fancy}
\fancyhf{}
\fancyhead[LE]{\small\textit{\nouppercase{\leftmark}}}   % chapter on even
\fancyhead[RO]{\small\textit{\nouppercase{\rightmark}}}  % section on odd
\fancyfoot[LE,RO]{\small\thepage}
\renewcommand{\headrulewidth}{0.4pt}

\fancypagestyle{plain}{
    \fancyhf{}
    \fancyfoot[C]{\small\thepage}
    \renewcommand{\headrulewidth}{0pt}
}

% ============================================================================
% DOCUMENT
% ============================================================================

\begin{document}

% %%%%%%%%%%%%%%%%%%%%%%%%%%%%%%%%%%%%%%%%%%%%%%%%%%%%%%%%%%%%%%%%%%%%%%%%%%%%
% FRONT MATTER
% Pages numbered with roman numerals (i, ii, iii, ...)
% %%%%%%%%%%%%%%%%%%%%%%%%%%%%%%%%%%%%%%%%%%%%%%%%%%%%%%%%%%%%%%%%%%%%%%%%%%%%

\frontmatter

% --------------------------------------------------------------------------
% Title Page
% --------------------------------------------------------------------------
\begin{titlepage}
\centering

\vspace*{2cm}

{\color{thesisblue}\rule{\textwidth}{1pt}}

\vspace{1.5cm}

{\LARGE\bfseries On the Convergence of Iterative Methods\\[0.4em]
for Sparse Linear Systems}

\vspace{1.5cm}

{\color{thesisblue}\rule{\textwidth}{1pt}}

\vspace{2cm}

{\large A thesis submitted in partial fulfilment of the requirements for\\[0.3em]
the degree of \textbf{Doctor of Philosophy}}

\vspace{2cm}

{\Large Jane Doe}

\vspace{1cm}

{\large Department of Mathematics\\[0.2em]
Example University\\[0.2em]
\today}

\vfill

{\small Supervisor: Prof.\ John Smith \quad Co-supervisor: Dr.\ Alice Brown}

\end{titlepage}

% --------------------------------------------------------------------------
% Declaration
% --------------------------------------------------------------------------
\chapter*{Declaration}
\addcontentsline{toc}{chapter}{Declaration}

I declare that this thesis is my own work and has not been submitted for any
other degree or professional qualification. All sources are acknowledged in
the bibliography.

\vspace{2cm}
\noindent
\rule{6cm}{0.4pt}\\
Jane Doe\\
\today

% --------------------------------------------------------------------------
% Abstract
% --------------------------------------------------------------------------
\chapter*{Abstract}
\addcontentsline{toc}{chapter}{Abstract}

This thesis investigates the convergence properties of preconditioned
Krylov-subspace methods applied to large sparse linear systems arising from
the discretisation of elliptic partial differential equations.  We derive
tight convergence bounds for the Conjugate Gradient method under a class of
incomplete Cholesky preconditioners, and propose a hierarchical coarsening
strategy that reduces the number of iterations by up to 40\% on a benchmark
set of 50 test matrices.

\medskip
\noindent\textbf{Keywords:} sparse linear systems, Krylov-subspace methods,
preconditioning, convergence analysis, incomplete Cholesky

% --------------------------------------------------------------------------
% Acknowledgements
% --------------------------------------------------------------------------
\chapter*{Acknowledgements}
\addcontentsline{toc}{chapter}{Acknowledgements}

I am deeply grateful to my supervisor Prof.\ John Smith for his guidance,
patience, and sharp mathematical insight throughout this project.  I thank
Dr.\ Alice Brown for many stimulating discussions, and the Numerical Analysis
Group for creating a wonderfully collaborative environment.

Financial support from the Example Science Foundation (grant ES-12345) is
gratefully acknowledged.

% --------------------------------------------------------------------------
% Table of contents, figures, tables
% --------------------------------------------------------------------------
\tableofcontents
\listoffigures
\listoftables

% %%%%%%%%%%%%%%%%%%%%%%%%%%%%%%%%%%%%%%%%%%%%%%%%%%%%%%%%%%%%%%%%%%%%%%%%%%%%
% MAIN MATTER
% Pages numbered with arabic numerals (1, 2, 3, ...)
% %%%%%%%%%%%%%%%%%%%%%%%%%%%%%%%%%%%%%%%%%%%%%%%%%%%%%%%%%%%%%%%%%%%%%%%%%%%%

\mainmatter

% --------------------------------------------------------------------------
% Chapter 1 — Introduction
% --------------------------------------------------------------------------

\chapter{Introduction}
\label{chap:intro}

Large sparse linear systems of the form
\begin{equation}
    \label{eq:linear_system}
    A\mathbf{x} = \mathbf{b}, \quad A \in \R^{n \times n}, \quad
    \mathbf{b} \in \R^{n},
\end{equation}
arise ubiquitously in the numerical solution of partial differential
equations (PDEs), circuit simulation, and structural mechanics.  Direct
methods based on Gaussian elimination are prohibitively expensive when $n$
exceeds~$10^6$; iterative methods that exploit sparsity are then the only
viable approach \cite{saad2003iterative}.

\section{Motivation}
\label{sec:intro_motivation}

The central challenge is to accelerate convergence by constructing a
preconditioner $M \approx A$ such that the preconditioned system
$M^{-1}A\mathbf{x} = M^{-1}\mathbf{b}$ has a favourable eigenvalue
distribution.  \Cref{chap:methods} reviews classical preconditioners;
\cref{chap:results} presents the new hierarchical approach.

\section{Contributions}
\label{sec:intro_contributions}

The main contributions of this thesis are:
\begin{enumerate}
    \item A unified convergence framework for the class of $\ell$-level
          incomplete Cholesky preconditioners (\cref{thm:convergence_bound}).
    \item A hierarchical coarsening algorithm (HC-IC) that adaptively selects
          the fill-in level based on local spectral estimates.
    \item A systematic numerical comparison on the SuiteSparse Matrix
          Collection, demonstrating up to 40\% iteration reduction.
\end{enumerate}

\section{Outline}
\label{sec:intro_outline}

The remainder of this thesis is organised as follows.
\Cref{chap:background} reviews the mathematical background.
\Cref{chap:methods} develops the new preconditioner.
\Cref{chap:results} reports numerical experiments.
\Cref{chap:conclusion} concludes and suggests future directions.

% --------------------------------------------------------------------------
% Chapter 2 — Background
% --------------------------------------------------------------------------

\chapter{Background}
\label{chap:background}

\section{Sparse Matrix Storage}
\label{sec:storage}

A matrix $A$ is \emph{sparse} if the number of non-zero entries
$\mathrm{nnz}(A) \ll n^2$.  Common storage formats include:

\begin{description}
    \item[CSR (Compressed Sparse Row)] Three arrays store values, column
        indices, and row pointers. Row operations are cache-efficient.
    \item[CSC (Compressed Sparse Column)] The column-oriented analogue of CSR.
        Preferred for column operations (e.g., in Fortran-ordered solvers).
    \item[COO (Coordinate)] Explicit $(i, j, v)$ triples; simple but slower
        for arithmetic.
\end{description}

\section{Krylov-Subspace Methods}
\label{sec:krylov}

Given an initial guess $\mathbf{x}_0$, the residual is
$\mathbf{r}_0 = \mathbf{b} - A\mathbf{x}_0$.  The $k$-th Krylov subspace is
\[
    \mathcal{K}_k(A, \mathbf{r}_0)
    = \mathrm{span}\{\mathbf{r}_0,\, A\mathbf{r}_0,\, \ldots,\,
                     A^{k-1}\mathbf{r}_0\}.
\]
The Conjugate Gradient (CG) method is applicable when $A$ is
symmetric positive definite (SPD); GMRES is used for non-symmetric systems.

\begin{theorem}[CG Convergence Bound]
    \label{thm:convergence_bound}
    Let $A$ be SPD with condition number $\kappa(A)$.  After $k$ CG
    iterations,
    \[
        \frac{\norm{\mathbf{e}_k}_A}{\norm{\mathbf{e}_0}_A}
        \leq 2\left(\frac{\sqrt{\kappa(A)}-1}{\sqrt{\kappa(A)}+1}\right)^k,
    \]
    where $\norm{\mathbf{v}}_A = \sqrt{\mathbf{v}^\top A\mathbf{v}}$ denotes
    the $A$-norm.
\end{theorem}

\section{Preconditioning}
\label{sec:preconditioning}

A preconditioner $M$ approximates $A$ while being cheap to apply.  The
\emph{incomplete LU (ILU)} factorisation \cite{meijerink1977iterative}
computes $A \approx LU$ by ignoring fill-in outside a prescribed sparsity
pattern.  The \emph{incomplete Cholesky (IC)} factorisation is the SPD
analogue.

% --------------------------------------------------------------------------
% Chapter 3 — Methods
% --------------------------------------------------------------------------

\chapter{Methods}
\label{chap:methods}

\section{Hierarchical Coarsening Algorithm}
\label{sec:hc_algorithm}

The proposed HC-IC algorithm proceeds in two phases:

\begin{enumerate}
    \item \textbf{Coarsening.}  Construct a sequence of coarser graphs
          $G_0 \supset G_1 \supset \cdots \supset G_L$ using a maximal
          independent set (MIS) criterion.
    \item \textbf{Factorisation.}  Apply IC(0) on $G_L$ and transfer the
          factorisation upward, adding fill-in guided by spectral estimates
          at each level.
\end{enumerate}

\cref{alg:hcic} gives the pseudocode.

\begin{figure}[h]
\centering
% Placeholder box for an algorithm/diagram:
\fbox{\parbox{0.8\textwidth}{\centering\vspace{1cm}
\textit{[Algorithm pseudocode or diagram would appear here.]}
\vspace{1cm}}}
\caption{Pseudocode for the HC-IC preconditioner construction.}
\label{alg:hcic}
\end{figure}

\section{Convergence Analysis}
\label{sec:convergence_analysis}

\begin{lemma}
    \label{lem:eigenvalue_clustering}
    For a matrix $A$ with hierarchically coarsened IC preconditioner $M$,
    the eigenvalues of $M^{-1}A$ cluster in the interval $[1-\epsilon, 1+\epsilon]$
    for a constant $\epsilon < 1$ that depends on the coarsening ratio.
\end{lemma}

\begin{proof}
    By induction on the coarsening level.  At level 0, $M = A$ and
    $\epsilon = 0$.  At each subsequent level, the additional fill-in
    reduces the residual factorisation error by a factor bounded by the
    spectral gap of the coarse grid operator.  The full proof follows
    the argument of \cite{vassilevski2008multilevel}.
\end{proof}

% --------------------------------------------------------------------------
% Chapter 4 — Numerical Results
% --------------------------------------------------------------------------

\chapter{Numerical Results}
\label{chap:results}

\section{Test Problems}
\label{sec:test_problems}

We selected 50 SPD matrices from the SuiteSparse Matrix Collection, spanning
dimensions from $10^4$ to $10^7$ and condition numbers from $10^2$ to $10^9$.
\Cref{tab:test_matrices} summarises three representative cases.

\begin{table}[h]
\centering
\caption{Representative test matrices from SuiteSparse.}
\label{tab:test_matrices}
\begin{tabular}{lrrrc}
\toprule
\textbf{Matrix} & \boldmath$n$ & \boldmath$\mathrm{nnz}(A)$ & \boldmath$\kappa(A)$ & \textbf{Origin} \\
\midrule
\texttt{thermal2}   & 1\,228\,045 & 8\,580\,313  & $1.2 \times 10^4$ & Heat transfer FEM \\
\texttt{ecology2}   & 999\,999    & 4\,995\,991  & $6.7 \times 10^5$ & Ecological diffusion \\
\texttt{G3\_circuit} & 1\,585\,478 & 7\,660\,826 & $3.8 \times 10^8$ & Circuit simulation \\
\bottomrule
\end{tabular}
\end{table}

\section{Performance Comparison}
\label{sec:performance}

\Cref{tab:iterations} compares iteration counts for IC(0), IC(1), and the
proposed HC-IC preconditioner with CG on the three test problems.

\begin{table}[h]
\centering
\caption{CG iteration counts with different preconditioners.
         Convergence tolerance: $\norm{\mathbf{r}_k}/\norm{\mathbf{r}_0} \leq 10^{-8}$.}
\label{tab:iterations}
\begin{tabular}{lrrr}
\toprule
\textbf{Matrix} & \textbf{IC(0)} & \textbf{IC(1)} & \textbf{HC-IC (ours)} \\
\midrule
\texttt{thermal2}    & 312 & 187 & \textbf{134} \\
\texttt{ecology2}    & 841 & 503 & \textbf{298} \\
\texttt{G3\_circuit} & --- & --- & \textbf{712} \\
\bottomrule
\multicolumn{4}{l}{\footnotesize --- = did not converge within 2000 iterations.}
\end{tabular}
\end{table}

\section{Discussion}
\label{sec:discussion}

HC-IC consistently reduces iteration counts compared with flat IC(0) and
IC(1), with the largest gains for ill-conditioned problems such as
\texttt{G3\_circuit}.  The overhead per iteration is approximately 15\%
higher than IC(0), but total solve time is still reduced because the
iteration count reduction dominates.

% --------------------------------------------------------------------------
% Chapter 5 — Conclusion
% --------------------------------------------------------------------------

\chapter{Conclusion}
\label{chap:conclusion}

\section{Summary}
\label{sec:summary}

This thesis has presented a hierarchical coarsening strategy for incomplete
Cholesky preconditioning of large sparse SPD systems.  The key results are:

\begin{itemize}
    \item A rigorous convergence framework (\cref{thm:convergence_bound,lem:eigenvalue_clustering})
    \item The HC-IC algorithm (\cref{sec:hc_algorithm}) that reduces CG
          iterations by 20--40\% on benchmark problems
    \item Open-source implementation available at
          \url{https://github.com/example/hcic}
\end{itemize}

\section{Future Work}
\label{sec:future}

Promising directions for future investigation include:
\begin{enumerate}
    \item Extension to non-symmetric systems using HC-ILU
    \item GPU acceleration of the coarsening and solve phases
    \item Adaptive tolerance selection based on online condition number
          estimates
\end{enumerate}

% %%%%%%%%%%%%%%%%%%%%%%%%%%%%%%%%%%%%%%%%%%%%%%%%%%%%%%%%%%%%%%%%%%%%%%%%%%%%
% BACK MATTER
% Page numbers continue from main matter
% %%%%%%%%%%%%%%%%%%%%%%%%%%%%%%%%%%%%%%%%%%%%%%%%%%%%%%%%%%%%%%%%%%%%%%%%%%%%

\backmatter

% --------------------------------------------------------------------------
% Appendices
% --------------------------------------------------------------------------

\appendix

\chapter{Proof of Lemma~\ref{lem:eigenvalue_clustering}}
\label{app:proof}

This appendix gives the full inductive proof of
\cref{lem:eigenvalue_clustering} omitted from the main text.

\section{Base Case}

At level $L$ (coarsest grid), IC is applied with zero fill-in.  The
eigenvalues of $M_L^{-1}A_L$ satisfy the claimed bound with
$\epsilon_L = \sqrt{1 - \lambda_{\min}(M_L^{-1}A_L)}$.

\section{Inductive Step}

Assume the bound holds at level $\ell+1$.  The level-$\ell$ preconditioner
is constructed by injecting fill-in based on spectral estimates from
level~$\ell+1$.  By the Weyl perturbation theorem, the eigenvalue
perturbation at level~$\ell$ satisfies $\delta\lambda \leq \rho \cdot
\epsilon_{\ell+1}$ for a contraction factor $\rho < 1$ that depends only
on the coarsening ratio.  Setting $\epsilon_\ell = \rho^{L-\ell}\epsilon_L$
completes the induction. \hfill$\square$

\chapter{Notation}
\label{app:notation}

\begin{center}
\begin{tabular}{ll}
\toprule
\textbf{Symbol} & \textbf{Meaning} \\
\midrule
$A \in \R^{n \times n}$     & Square real coefficient matrix \\
$\mathbf{x}, \mathbf{b} \in \R^n$ & Solution and right-hand side vectors \\
$\kappa(A) = \sigma_{\max}/\sigma_{\min}$ & Spectral condition number \\
$\norm{\cdot}$              & Euclidean ($\ell_2$) norm unless noted \\
$\norm{\cdot}_A$            & $A$-energy norm: $\norm{\mathbf{v}}_A = \sqrt{\mathbf{v}^\top A \mathbf{v}}$ \\
$\mathrm{nnz}(A)$           & Number of non-zero entries in $A$ \\
$\mathcal{K}_k(A,\mathbf{r})$ & Krylov subspace of degree $k$ \\
\bottomrule
\end{tabular}
\end{center}

% --------------------------------------------------------------------------
% Bibliography
% --------------------------------------------------------------------------

\bibliographystyle{plain}
\bibliography{thesis_refs}

% If you do not have a .bib file, replace the two lines above with:
%
% \begin{thebibliography}{9}
%
% \bibitem{saad2003iterative}
% Yousef Saad.
% \textit{Iterative Methods for Sparse Linear Systems}.
% SIAM, 2nd edition, 2003.
%
% \bibitem{meijerink1977iterative}
% J.~A. Meijerink and H.~A. van~der Vorst.
% An iterative solution method for linear equation systems of which the
% coefficient matrix is a symmetric M-matrix.
% \textit{Mathematics of Computation}, 31(137):148--162, 1977.
%
% \bibitem{vassilevski2008multilevel}
% Panayot~S. Vassilevski.
% \textit{Multilevel Block Factorisation Preconditioners}.
% Springer, 2008.
%
% \end{thebibliography}

\end{document}

% ============================================================================
% COMPILATION NOTES
% ============================================================================
% Full sequence (with BibTeX):
%   pdflatex thesis_template.tex
%   bibtex   thesis_template
%   pdflatex thesis_template.tex
%   pdflatex thesis_template.tex
%
% With latexmk (recommended):
%   latexmk -pdf thesis_template.tex
%
% If thesis_refs.bib does not exist, use the inline thebibliography block
% in the bibliography section above instead.
% ============================================================================
